% !TeX root = ../main.tex
\chapter*{DANH MỤC TỪ VIẾT TẮT VÀ THUẬT NGỮ}
\addcontentsline{toc}{chapter}{Danh mục từ viết tắt và thuật ngữ}

\begin{longtable}{L{2.6cm} L{4.6cm} L{6.5cm}}
\toprule
\textbf{Viết tắt} & \textbf{Tiếng Anh đầy đủ} & \textbf{Giải thích} \\
\midrule
\endfirsthead
\toprule
\textbf{Viết tắt} & \textbf{Tiếng Anh đầy đủ} & \textbf{Giải thích} \\
\midrule
\endhead
\bottomrule
\endfoot

AI      & Artificial Intelligence            & Trí tuệ nhân tạo \\
CAM     & Class Activation Map               & Bản đồ kích hoạt theo lớp \\
CC      & Connected Component                & Thành phần liên thông \\
CNN     & Convolutional Neural Network       & Mạng nơ-ron tích chập \\
CRF     & Conditional Random Field           & Trường ngẫu nhiên có điều kiện \\
EDA     & Exploratory Data Analysis          & Phân tích khám phá dữ liệu \\
FN      & False Negative                     & Âm tính giả --- bỏ sót vùng bệnh \\
FP      & False Positive                     & Dương tính giả --- báo nhầm vùng lành \\
GPU     & Graphics Processing Unit           & Bộ xử lý đồ hoạ \\
Grad-CAM & Gradient-weighted Class Activation Mapping & Bản đồ kích hoạt có trọng số gradient \\
HSV     & Hue--Saturation--Value             & Không gian màu sắc độ--bão hoà--độ sáng \\
IoU     & Intersection over Union            & Tỉ số giao trên hợp \\
KDE     & Kernel Density Estimation          & Ước lượng mật độ hạt nhân \\
LR      & Learning Rate                      & Tốc độ học \\
ReLU    & Rectified Linear Unit              & Hàm kích hoạt tuyến tính chỉnh lưu \\
RGB     & Red--Green--Blue                   & Không gian màu đỏ--lục--lam \\
SAM     & Segment Anything Model             & Mô hình nền tảng phân đoạn tổng quát \\
TP      & True Positive                      & Dương tính thật \\
ViT     & Vision Transformer                 & Mạng Transformer cho thị giác máy tính \\
XAI     & Explainable Artificial Intelligence & Trí tuệ nhân tạo khả diễn giải \\
\end{longtable}

\vspace{0.6cm}
\noindent\textbf{Quy ước thuật ngữ.} Các thuật ngữ tiếng Anh đã trở thành chuẩn
trong lĩnh vực (\en{pseudo-label}, \en{fine-tuning}, \en{early stopping},
\en{ground truth}, \en{backbone}, \en{encoder}, \en{decoder}\dots) được giữ
nguyên dạng tiếng Anh và in nghiêng ở lần xuất hiện đầu tiên, kèm giải thích
tiếng Việt trong ngoặc. Ở các lần sau, thuật ngữ được dùng trực tiếp không in
nghiêng nhằm giữ mạch đọc.

\vspace{0.4cm}
\noindent\textbf{Quy ước ký hiệu phiên bản.} Luận văn xây dựng ba phiên bản
pipeline phân đoạn, ký hiệu thống nhất là \textbf{V1}, \textbf{V2}, \textbf{V3}.
Mỗi phiên bản gồm một bộ \en{pseudo-label} và một mô hình phân đoạn huấn luyện
trên bộ nhãn đó; chi tiết trong \cref{chap:phuongphap}.

\vspace{0.4cm}
\noindent\textbf{Quy ước số thập phân.} Luận văn dùng dấu phẩy làm dấu thập phân
theo chuẩn tiếng Việt (ví dụ: 0,6239). Trong các đoạn mã nguồn và tên tệp trích
dẫn nguyên trạng, dấu chấm được giữ nguyên theo đúng dữ liệu gốc.
